\chapter{Discussion and Recommendations}
\label{chap:discussion}

This chapter interprets the empirical findings presented in Chapter~3
in the context of the theoretical framework and existing literature
reviewed in Chapter~1. It examines the significance of the results,
compares them with prior research, proposes practical recommendations
for improving information system integration maturity in PMO environments
in Latvian enterprises, and reflects on unexpected findings and
theoretical implications.

\section{Discussion of Findings}
\label{sec:discussion}

\subsection{Integration Maturity Distribution in Latvian PMOs}
According to the survey findings, 70\% of the sampled PMO environments
operate at intermediate integration maturity levels --- either Level 2
(Developing) or Level 3 (Defined). This result supports Hypothesis 1
and is consistent with broader research on digital maturity in Baltic
enterprises, which indicates that Latvian organizations generally fall
behind their Nordic counterparts in terms of digital integration and
automation \cite{eurostat2023}.

The prevalence of Levels 2 and 3 indicates that most Latvian enterprises
have progressed beyond entirely manual data consolidation, yet full
platform-level integration has not been achieved. This is a crucial
observation, as it reveals that most organizations are in a transitional
phase --- they have acknowledged the need for integration but have not yet
committed fully to reaching Level 4. This result also aligns with the
expert interview, in which Expert 1 noted that challenges such as data
latency and legacy system resistance persist even in organizations with
relatively high integration maturity.

\subsection{Integration Maturity and Reporting Automation}
The statistically significant positive correlation between integration
maturity and reporting automation (r=0.405, p=0.010) confirms that
organizations with more advanced IS integration produce reports more
automatically. This supports Hypothesis 2 and is consistent with
research by Hohpe and Woolf \cite{hohpe2004}, who demonstrated that
API-based integration architectures substantially reduce the need for
manual data handling.

The practical implication is straightforward: organizations that invest
in higher IS integration levels can expect a notable reduction in the
time and effort devoted to manual reporting. Expert 1 confirmed this,
describing a 70/30 split in favour of automated reporting in their
organization, compared to what is likely a considerably lower automation
rate in Level 1 and Level 2 organizations.

\subsection{Integration Maturity and Data Quality}
The positive correlation between integration maturity and data quality
(r=0.321, p=0.043), while the weakest of the three, remains statistically
significant. This confirms Hypothesis 3 and supports the theoretical
argument presented in Chapter~1 that manual data consolidation introduces
errors and delays that compromise data quality \cite{dama2017}.

The comparatively weaker correlation for data quality may reflect the
fact that it is shaped by many factors beyond IS integration alone ---
including user behaviour, data entry discipline, and organizational
culture. Expert 1 emphasized this point, noting that even at high
integration maturity, data quality depends heavily on how consistently
project managers update their records. This suggests that technical
integration must be accompanied by organizational practices that
actively support data hygiene.

\subsection{Integration Maturity and Governance Standardization}
The strong positive correlation between integration maturity and
governance standardization (r=0.403, p=0.010) confirms Hypothesis 4
and provides empirical support for the theoretical link between
centralized BI systems and governance quality established in
Chapter~1 \cite{negash2004}. Organizations with higher integration
maturity demonstrated noticeably stronger governance standardization
practices.

The expert interview strongly supports this finding. Expert 1 described
how the introduction of a centralized Power BI Command Center transformed
governance in their organization by enforcing uniform KPI definitions,
eliminating subjective reporting, and making project performance
transparently visible to all executives at the same time. This
qualitative evidence helps explain the mechanism through which higher
integration maturity leads to improved governance standardization.

\section{Comparison with Existing Literature}
\label{sec:literature-comparison}

The findings of this study are broadly consistent with existing research
on IS integration and PMO maturity. The dominance of intermediate maturity
levels in the sample aligns with research by Hobbs and Aubry
\cite{hobbs2007}, who found that most PMOs struggle to advance beyond
basic standardization practices. The positive relationships between
integration maturity and the three dependent variables are consistent
with the theoretical predictions of the Integration Maturity Model
proposed in Chapter~1.

Nevertheless, this study makes a distinctive contribution by providing
empirical evidence specifically from the Latvian enterprise context.
As noted in Chapter~1, no prior study had examined IS integration
maturity in PMO environments within Latvia. The findings suggest that
Latvian enterprises follow a pattern similar to those documented in
broader European research, but with a tendency toward lower overall
maturity levels, which is consistent with the digital maturity gap
between Baltic and Nordic countries identified by Eurostat
\cite{eurostat2023}.

\section{Unexpected Findings and Theoretical Implications}
\label{sec:implications}

\subsection{Unexpected Findings}
A notable finding was that data quality showed the weakest correlation
with integration maturity (r=0.321) compared to reporting automation
(r=0.405) and governance standardization (r=0.403). This was somewhat
unexpected, as the theoretical framework suggested that automated
integration would directly improve data accuracy by eliminating manual
entry errors. However, as Expert 1 observed, data quality is also
strongly shaped by human behaviour and organizational discipline ---
factors that operate independently of technical integration level.
This suggests that the relationship between IS integration and data
quality is more complex than initially anticipated and warrants
further investigation.

Another interesting observation is that 15\% of organizations operated
at Level 4 --- the highest maturity level. This proportion is higher than
might be expected given the generally lower digital maturity of Baltic
enterprises relative to Nordic counterparts. This may reflect a degree
of self-selection bias, as organizations with more advanced PMO
structures may have been more willing to participate in a PMO-focused
survey.

\subsection{Theoretical Implications}
This study advances the theoretical understanding of IS integration
maturity in PMO environments in several ways. First, it provides
empirical validation of the four-level Integration Maturity Model
proposed in Chapter~1, demonstrating that its levels correspond to
meaningful differences in reporting automation, data quality, and
governance standardization. Second, it extends the existing body of
literature on digital maturity models by applying the maturity concept
specifically to PMO IS integration --- a context that had not previously
been examined empirically. Third, it contributes the first empirical
evidence of IS integration maturity patterns in Latvian enterprises,
adding to the growing body of research on digital transformation in
Baltic and Eastern European contexts.

\section{Practical Recommendations}
\label{sec:recommendations}

Based on the empirical findings and expert interview insights, the
following practical recommendations are proposed for medium and large
Latvian enterprises seeking to improve their PMO IS integration maturity.

\subsection{Recommendation 1 --- Prioritize API-Based Integration}
Organizations currently at Level 1 or Level 2 should focus on
establishing API-based connections between their core PMO systems ---
particularly between project tracking tools, financial systems, and
reporting platforms. Even basic API integration can substantially
reduce manual data consolidation and improve reporting automation,
as supported by the strong correlation between maturity level and
reporting automation found in this study \cite{hohpe2004}.

\subsection{Recommendation 2 --- Invest in Centralized BI Dashboards}
One of the most effective steps organizations can take to strengthen
governance standardization is to adopt a centralized BI or dashboard
solution. Tools such as Power BI, Tableau, or Smartsheet can bring
together data from multiple systems into a unified view, enforcing
consistent KPI definitions and improving transparency across the
organization \cite{negash2004}. As Expert 1 noted, this single change
can be transformative for governance quality.

\subsection{Recommendation 3 --- Consider Low-Code Integration Platforms}
For medium-sized Latvian enterprises that lack the resources for
large-scale enterprise integration, low-code and modular platforms
offer a cost-effective route toward higher integration maturity.
During the expert interview, Expert 1 specifically recommended
platforms such as Monday.com \cite{monday2024}, ClickUp
\cite{clickup2024}, and Zapier \cite{zapier2024} as practical
options that enable organizations to connect existing systems without
requiring extensive technical expertise or significant financial
investment. These platforms are particularly well suited to
organizations at Level 2 seeking a pragmatic path toward Level 3
integration.

\subsection{Recommendation 4 --- Address Local System Interoperability}
Many Latvian enterprises rely on local accounting and ERP software
that is not natively compatible with global PM tools. Organizations
should prioritize developing or procuring standardized connectors
between these local systems and their PMO tools to eliminate the
manual data exports that currently fill this gap. Expert 1 highlighted
this as one of the most persistent practical barriers to integration
maturity advancement in the Latvian enterprise context.

\subsection{Recommendation 5 --- Invest in Data Literacy}
Technical integration alone is insufficient to improve data quality.
Organizations must also invest in training project managers and PMO
staff to appreciate the importance of data hygiene and to use
integrated systems effectively \cite{dama2017}. As Expert 1 observed,
the introduction of a Data Health Score system helped sustain data
accuracy even in a highly integrated environment, demonstrating that
organizational practices must accompany technical solutions.

\section{Limitations of the Study}
\label{sec:study-limitations}

Although the findings of this study are meaningful and statistically
significant, several limitations should be kept in mind when
interpreting the results. The sample of 40 respondents is adequate
for the planned analyses but restricts the statistical power of more
advanced models. Reliance on self-reported data means that respondents
may have overestimated their organization's integration maturity. The
cross-sectional design captures only a single point in time and cannot
establish causal relationships between variables. Finally, the
dominance of IT sector organizations in the sample may limit the
generalizability of findings to other industries. These limitations
are discussed in greater detail in Chapter~2.

\section{Chapter Summary}
\label{sec:chap4summary}

This chapter examined the empirical findings of the study within the
context of the theoretical framework and existing literature. All four
hypotheses were confirmed by the data, with statistically significant
positive correlations identified between IS integration maturity and
reporting automation, data quality, and governance standardization.
The findings align with broader research on PMO maturity and digital
transformation in Baltic enterprises. A noteworthy result was the
comparatively weak correlation between integration maturity and data
quality, suggesting that human behaviour and organizational culture
play an important role alongside technical integration. Five practical
recommendations were developed to support Latvian enterprises in
advancing their PMO IS integration maturity, with specific tools
identified based on both the empirical findings and the qualitative
insights provided by Expert 1. These findings and recommendations
provide the foundation for the conclusions presented in the following
chapter.
